% ===========================================================================
%  SCX-Audited Quantitative Finance:
%  Multi-Model Consensus for Derivative Pricing and Risk Certification
%  —— MATH-DRIVEN, not survey. Theorems + proofs + explicit assumptions.
%  SCX, June 2026
% ===========================================================================

\documentclass[11pt,a4paper]{article}

% ─── Packages ───────────────────────────────────────────────────────────
\usepackage[utf8]{inputenc}
\usepackage[T1]{fontenc}
\usepackage{amsmath,amssymb,amsthm,mathtools}
\usepackage{bm}
\usepackage{booktabs}
\usepackage{graphicx}
\usepackage{hyperref}
\usepackage[margin=2.5cm]{geometry}
\usepackage{enumitem}
\usepackage{cleveref}
\usepackage{xcolor}
\usepackage{setspace}
\usepackage{multirow}
\usepackage{array}
\usepackage{algorithm}
\usepackage{algpseudocode}
\usepackage{stmaryrd}
\usepackage{dsfont}
\onehalfspacing

% ─── Theorem Environments ──────────────────────────────────────────────
\theoremstyle{plain}
\newtheorem{theorem}{Theorem}[section]
\newtheorem{lemma}[theorem]{Lemma}
\newtheorem{corollary}[theorem]{Corollary}
\newtheorem{proposition}[theorem]{Proposition}
\newtheorem{conjecture}[theorem]{Conjecture}

\theoremstyle{definition}
\newtheorem{definition}[theorem]{Definition}
\newtheorem{assumption}[theorem]{Assumption}

\theoremstyle{remark}
\newtheorem{remark}[theorem]{Remark}
\newtheorem{example}[theorem]{Example}
\newtheorem{openproblem}[theorem]{Open Problem}

% ─── Rigor Status Labels ────────────────────────────────────────────────
\definecolor{rigorFull}{RGB}{0,100,0}
\definecolor{rigorPartial}{RGB}{180,120,0}
\definecolor{rigorSketch}{RGB}{150,0,0}
\newcommand{\rigorFull}{{\color{rigorFull}$\blacksquare$\,\textbf{[Rigor: FULL $\star\star\star$]}}}
\newcommand{\rigorPartial}{{\color{rigorPartial}$\blacksquare$\,\textbf{[Rigor: PARTIAL $\star\star\;$]}}}
\newcommand{\rigorSketch}{{\color{rigorSketch}$\blacksquare$\,\textbf{[Rigor: SKETCH $\star\;\,$]}}}

% ─── Math Operators ─────────────────────────────────────────────────────
\DeclareMathOperator*{\argmin}{arg\,min}
\DeclareMathOperator*{\argmax}{arg\,max}
\DeclareMathOperator{\TV}{TV}
\DeclareMathOperator{\KL}{D_{\mathrm{KL}}}
\DeclareMathOperator{\Var}{Var}
\DeclareMathOperator{\Bias}{Bias}
\DeclareMathOperator{\Cov}{Cov}
\DeclareMathOperator{\RMSE}{RMSE}
\DeclareMathOperator{\Corr}{Corr}
\DeclareMathOperator{\sgn}{sgn}
\DeclareMathOperator{\diag}{diag}
\DeclareMathOperator{\Lip}{Lip}

% ─── Custom Commands ────────────────────────────────────────────────────
\newcommand{\R}{\mathbb{R}}
\newcommand{\N}{\mathbb{N}}
\newcommand{\E}{\mathbb{E}}
\newcommand{\Pbb}{\mathbb{P}}
\newcommand{\F}{\mathcal{F}}
\newcommand{\X}{\mathcal{X}}
\newcommand{\Y}{\mathcal{Y}}
\newcommand{\Sstates}{\mathcal{S}}
\newcommand{\D}{\mathcal{D}}
\newcommand{\M}{\mathcal{M}}
\newcommand{\ind}[1]{\mathbf{1}\{#1\}}
\newcommand{\norm}[1]{\|#1\|}
\newcommand{\abs}[1]{\lvert#1\rvert}
\newcommand{\inner}[2]{\langle #1, #2 \rangle}
\newcommand{\Situs}{\textsc{Situs}}
\newcommand{\Yajie}{\textsc{Yajie}}
\newcommand{\Spring}{\textsc{Spring}}
\newcommand{\Cercis}{\textsc{Cercis}}
\newcommand{\SCX}{\textsc{SCX}}

% Finance-specific operators
\DeclareMathOperator{\BS}{BS}
\DeclareMathOperator{\Heston}{Heston}
\DeclareMathOperator{\SABR}{SABR}
\DeclareMathOperator{\rBergomi}{rBergomi}
\newcommand{\Crank}{\mathrm{C}_{\mathrm{rank}}}
\newcommand{\Pmiss}{\mathrm{P}_{\mathrm{miss}}}
\newcommand{\Dmis}{\Delta_{\mathrm{mis}}}
\newcommand{\sigmamkt}{\sigma_{\mathrm{mkt}}}
\newcommand{\bidask}{s_{\mathrm{BA}}}
\newcommand{\hestG}{\boldsymbol{\Theta}_{\mathrm{H}}}
\newcommand{\sabrG}{\boldsymbol{\Theta}_{\mathrm{S}}}
\newcommand{\svjjG}{\boldsymbol{\Theta}_{\mathrm{SVJJ}}}
\newcommand{\rbergomiG}{\boldsymbol{\Theta}_{\mathrm{rB}}}

% ─── Reference to SCX Core Theorems ─────────────────────────────────────
\newcommand{\ThmSCXNoise}{Theorem~1 (SCX Noise Detection)}
\newcommand{\ThmSCXWeakFeature}{Theorem~2 (SCX Weak Features)}
\newcommand{\ThmSCXHonest}{Theorem~3 (SCX Unidentifiability)}
\newcommand{\ThmSCXPole}{Theorem~4 (SCX Pole)}

% ===========================================================================
% TITLE
% ===========================================================================
\title{\textbf{SCX-Audited Quantitative Finance: Multi-Model Consensus\\
       for Derivative Pricing and Risk Certification}\\[4pt]
       {\large —— SCX审计量化金融：衍生品定价与风险认证的多模型共识 ——}}

\author{SCX Quantitative Finance Division\\
        Nous Research --- SCX Audit Division\\[4pt]
        \texttt{scx-quant@nousresearch.com}}
\date{June 30, 2026}

\begin{document}
\maketitle

% ===========================================================================
% ABSTRACT
% ===========================================================================
\begin{abstract}
\noindent\textbf{English.}
Derivative pricing models---Black-Scholes, Heston stochastic volatility, SABR,
SVJJ with jumps, and rough Bergomi fractional volatility---each capture distinct
features of market dynamics, yet none dominates across all market regimes. When
these models disagree, traders and risk managers face an auditable decision
problem: is the disagreement a genuine mispricing opportunity, a consequence of
model misspecification, or a signal of an unprecedented structural regime shift?
We formalize derivative pricing as a multi-expert prediction problem under the
SCX (Structured Causal eXamination) framework and prove three fundamental
theorems. \textbf{Theorem~1} (Multi-Model Mispricing Detection): for $M$
independently-calibrated pricing models with disagreement exceeding the observed
bid-ask spread, the probability of missing a genuine arbitrage-adjacent
mispricing is bounded by the Hoeffding concentration
$P(\text{miss}) \leq \exp(-2M \cdot (\Delta / \sigma_{\mathrm{mkt}})^2)$,
where $\Delta$ is the inter-model disagreement and $\sigma_{\mathrm{mkt}}$ is
the market-implied volatility scale. \textbf{Theorem~2} (Model Risk vs.\ Regime
Shift Unidentifiability): when {\em all} $M$ models fail systematically---as
occurred in March 2020 for SPX options---the source of failure (collective model
misspecification versus a genuine structural break in the volatility process) is
provably unidentifiable from price data alone, providing a formal connection to
the Lucas Critique in macroeconomics. \textbf{Theorem~3} (Spring Gating for
Volatility Regime Detection): the \Spring{} gating mechanism detects volatility
regime transitions through a trajectory-divergence statistic on the
multi-expert prediction manifold, with false-alarm rate bounded by Hoeffding's
inequality on the empirical Wasserstein distance. We define the \Cercis{} Score
for quantitative finance as $S = Q + \eta N$, where $Q$ is a composite of
calibration error and hedging error across all models, and $N$ is a
regime-novelty metric based on the maximum mean discrepancy between current and
historical market states. Empirical benchmarks on SPX options (2004--2024),
VIX derivatives, interest rate swaptions, and single-name CDS spreads
demonstrate that SCX-audited multi-model consensus provides calibrated
mispricing detection and early regime-shift warning that single-model approaches
cannot achieve.

\medskip
\noindent\textbf{Chinese Abstract (中文摘要).}
衍生品定价模型——Black-Scholes、Heston随机波动率、SABR、带跳跃的SVJJ和粗糙Bergomi分形波动率——各自捕捉了市场动力学的不同特征，但没有任何模型在所有市场状态下都占主导地位。当这些模型出现分歧时，交易员和风险管理者面临一个可审计的决策问题：分歧是真正的错误定价机会、模型设定错误的后果，还是前所未有的结构性状态转换信号？我们在SCX框架下将衍生品定价形式化为多专家预测问题，证明了三个基本定理。\textbf{定理1（多模型错误定价检测）}：对$M$个独立校准的定价模型，当模型间分歧超过买卖价差时，遗漏真实套利相邻错误定价的概率受Hoeffding集中不等式约束。\textbf{定理2（模型风险与状态转换不可辨识性）}：当所有$M$个模型系统性地失败时（如2020年3月SPX期权），失败来源——集体模型误设还是波动率过程的真正结构性突变——从价格数据本身是无法辨识的，这为宏观经济学的卢卡斯批判提供了形式化联系。\textbf{定理3（Spring门控波动率状态检测）}：通过多专家预测流形上的轨迹分歧统计量检测波动率状态转换。我们定义了量化金融的\Cercis{}评分$S=Q+\eta N$，其中$Q$包含所有模型的校准误差和对冲误差，$N$是基于当前与历史市场状态最大均值差异的状态新奇度量。在SPX期权（2004--2024）、VIX衍生品、利率互换期权和单名CDS价差的实证基准测试表明，SCX审计的多模型共识提供了单一模型方法无法实现的可校准错误定价检测和早期状态转换预警。

\medskip
\noindent\textbf{Keywords:} 量化金融 (quantitative finance), 衍生品定价 (derivative pricing),
模型风险 (model risk), 多模型共识 (multi-model consensus), 卢卡斯批判 (Lucas Critique),
SCX审计 (SCX audit), 波动率状态检测 (volatility regime detection),
Cercis评分 (Cercis Score), 弹簧门控 (Spring gating), 买卖价差 (bid-ask spread).
\end{abstract}

% ======================================================================
\section{Introduction 引言}
\label{sec:intro}
% ======================================================================

Derivative pricing is the central computational task of quantitative finance.
Options, swaptions, credit default swaps, and exotic structured products are
priced daily by banks, hedge funds, and market makers operating in
multi-trillion-dollar markets. The canonical models---Black-Scholes (1973),
Heston stochastic volatility (1993), SABR (2002), SVJJ with jumps in returns and
volatility (Duffie, Pan, and Singleton, 2000), and rough Bergomi fractional
volatility (Gatheral, Jaisson, and Rosenbaum, 2018)---each capture different
aspects of the empirical volatility surface: the skew, the term structure of
kurtosis, the short-maturity explosion of the at-the-money skew, and the
clustering of extreme events \cite{black1973,heston1993,hagan2002,duffie2000,
gatheral2018}.

Yet a fundamental problem persists: {\em no single model is correct across all
market regimes.} The Black-Scholes model, despite its analytic elegance, fails
to reproduce the volatility smile. The Heston model captures the smile but not
the short-maturity skew explosion that rough Bergomi models explain. SVJJ models
fit the joint dynamics of SPX and VIX but involve complex jump structures that
are difficult to calibrate robustly. Traders routinely use multiple models
simultaneously, manually comparing their outputs and applying ``trader
adjustments'' based on experience and market color.

This paper addresses three questions:

\begin{enumerate}[label=(\roman*)]
    \item \textbf{Mispricing detection (错误定价检测):} When $M$ independently-calibrated
    pricing models disagree on the fair value of a derivative beyond the
    observed bid-ask spread, what is the probability that this disagreement
    represents a genuine mispricing rather than model noise? (Theorem~1)

    \item \textbf{Regime-shift unidentifiability (状态转换不可辨识性):} When all $M$ models
    fail simultaneously---as occurred in March 2020 for SPX options, where
    every model under-predicted the implied volatility spike---can we distinguish
    collective model misspecification from a genuine structural break in the
    volatility process? (Theorem~2)

    \item \textbf{Regime detection (状态检测):} Can we construct a principled, non-parametric
    mechanism for detecting transitions between volatility regimes (e.g., from
    low-volatility ``normal'' markets to high-volatility crisis regimes) using
    only the trajectory of multi-model predictions? (Theorem~3)
\end{enumerate}

We answer these questions within the \SCX{} (Structured Causal eXamination)
framework \cite{scx2025}, which provides a theorem-driven methodology for
multi-expert auditing. SCX decomposes any prediction problem into: (i) state
discovery via \Situs{} encoding; (ii) multi-expert consensus via \Yajie{}
voting; (iii) regime gating via \Spring{} divergence detection; and
(iv) quality scoring via the \Cercis{} composite metric.

The core insight is that pricing models function as {\em domain experts} in
the SCX sense: they are independently calibrated predictive systems whose
outputs (option prices, Greeks, implied volatilities) provide a multi-view
representation of the same underlying market state. When models disagree
beyond the market's own friction (the bid-ask spread), the disagreement carries
information. When they all fail together, the failure reveals a fundamental
limit on what models can know.

\textbf{Paper structure.} Section~\ref{sec:related} reviews related work.
Section~\ref{sec:framework} formalizes the quantitative finance state space and
the multi-model expert architecture. Section~\ref{sec:assumptions} states the
eight core assumptions. Section~\ref{sec:theorems} presents the three main
theorems with complete proofs. Section~\ref{sec:cercis} defines the \Cercis{}
Score for quantitative finance. Section~\ref{sec:yajie} describes the \Yajie{}
consensus mechanism for model aggregation. Section~\ref{sec:spring} develops
the \Spring{} gating detector. Section~\ref{sec:situs} specifies the \Situs{}
encoding for financial instruments. Section~\ref{sec:benchmarks} reports
empirical benchmarks. Section~\ref{sec:discussion} concludes.

% ======================================================================
\section{Related Literature 相关文献}
\label{sec:related}
% ======================================================================

Our work bridges several strands of quantitative finance and machine learning.

\textbf{Option pricing models.} The Black-Scholes model \cite{black1973}
assumes constant volatility and continuous trading, yielding the celebrated
closed-form formula. The Heston model \cite{heston1993} introduces stochastic
volatility following a Cox-Ingersoll-Ross (CIR) process, capturing the
volatility smile. The SABR model \cite{hagan2002} provides asymptotic
approximations for stochastic alpha-beta-rho dynamics, widely used in interest
rate derivatives. SVJJ models \cite{duffie2000} add compound Poisson jumps to
both returns and volatility, explaining the joint dynamics of the SPX index
and VIX. Rough Bergomi \cite{gatheral2018} models volatility as a fractional
Brownian motion with Hurst parameter $H < 1/2$, explaining the short-maturity
explosion of the at-the-money skew. Bacry, Muzy, and coworkers \cite{bacry2015}
develop Hawkes process models that capture self-exciting behavior in financial
markets.

\textbf{Model risk.} The Basel Committee on Banking Supervision
\cite{basel2016} defines model risk as ``the potential for adverse consequences
from decisions based on incorrect or misused model outputs.'' Cont \cite{cont2006}
proposes a quantitative framework for model risk measurement via worst-case
scenario analysis. Derman \cite{derman1996} provides a practitioner's taxonomy
of model risk sources. Glasserman and Xu \cite{glasserman2014} develop
robustness bounds for risk measures under model uncertainty. Our work extends
this literature by providing {\em finite-sample guarantees} on the detection
of model disagreement (Theorem~1) and proving a fundamental {\em negative result}
on the identifiability of model failure sources (Theorem~2).

\textbf{Model combination in finance.} Bayesian model averaging \cite{hoeting1999}
provides a probabilistic framework for combining multiple models, but requires
specification of prior model probabilities and assumes the true model is in the
candidate set. Our \Yajie{} consensus mechanism does not require the true model
to be in the set---it provides an audit, not a blend. The forecast combination
literature \cite{timmermann2006} establishes that equal-weighted combinations
often outperform any single forecaster; we extend this with exponential
concentration guarantees for mispricing detection.

\textbf{Regime-switching models.} Markov-switching models \cite{hamilton1989}
allow the data-generating process to transition between discrete states with
constant transition probabilities. Smooth-transition models \cite{terasvirta1994}
allow continuous transitions. Our \Spring{} gating mechanism differs in that it
detects regime transitions without specifying the number or nature of regimes
a priori---it uses the multi-model prediction trajectory as the detection
signal.

\textbf{Lucas Critique.} The Lucas Critique \cite{lucas1976} established that
parameters of reduced-form econometric models change when policy regimes change.
We prove a formal analog for quantitative finance (Theorem~2): when the
volatility process undergoes a structural break, the parameters of all pricing
models shift, and the shift is unidentifiable from collective model
misspecification. See also the SCX economics paper \cite{scx_economics2026}
for a parallel development in macroeconomic forecasting.

\textbf{Information geometry in finance.} Amari and Nagaoka \cite{amari2000}
develop information geometry, providing a Riemannian structure on statistical
manifolds. Brigo and coworkers \cite{brigo2019} apply differential geometry to
interest rate models. We use information-geometric concepts to define the
\Spring{} trajectory divergence statistic (Theorem~3).

% ======================================================================
\section{Framework and Notation 框架与符号}
\label{sec:framework}
% ======================================================================

\subsection{Market State Space 市场状态空间}

Let time be indexed by $t \in [0,T]$ on a filtered probability space
$(\Omega, \F, \{\F_t\}_{t \geq 0}, \Pbb)$ supporting all stochastic processes.
The market state at time $t$ is a vector:

\begin{equation}\label{eq:market_state}
    \mathbf{x}_t = \bigl(S_t, \sigma_t^{\mathrm{ATM}}, \mathcal{V}_t, r_t, q_t,
    \mathbf{w}_t\bigr) \in \X \subset \R^{d_x},
\end{equation}

where:
\begin{itemize}
    \item $S_t$ is the underlying asset price (e.g., SPX index level);
    \item $\sigma_t^{\mathrm{ATM}}$ is the at-the-money implied volatility at
    a benchmark maturity (e.g., 30-day);
    \item $\mathcal{V}_t = \{(\tau_i, \delta_j, \sigma_{\mathrm{imp}}(t; \tau_i,
    \delta_j))\}$ is the full implied volatility surface at time $t$, with
    $\tau_i$ the time-to-maturity and $\delta_j$ the option delta;
    \item $r_t$ is the risk-free interest rate;
    \item $q_t$ is the dividend yield;
    \item $\mathbf{w}_t \in \R^{d_w}$ is a vector of macro-financial covariates:
    VIX level, VIX term structure, credit spreads (CDX IG/HY), swaption implied
    volatility, Treasury term premium, and liquidity measures.
\end{itemize}

The full history up to time $t$ is denoted $\mathbf{x}_{[0,t]}$.

\subsection{Regime Definition 状态定义}

The market operates in one of $K$ discrete volatility regimes:
\begin{equation}\label{eq:regimes}
    \mathcal{R} = \{R_1, R_2, R_3, R_4, R_5\} = \{\text{Low Vol (低波动)},
    \text{Normal (正常)}, \text{Elevated (升高)}, \text{Crisis (危机)},
    \text{Tail (尾部)}\}.
\end{equation}

Regime membership is not directly observable but inferred from the implied
volatility surface $\mathcal{V}_t$ and macro-financial covariates
$\mathbf{w}_t$. The regime governs the data-generating process for asset
returns and volatility dynamics.

\subsection{Pricing Models as Experts 定价模型作为专家}

We deploy $M = 5$ option pricing models (``experts''):

\begin{equation}\label{eq:experts}
    \mathcal{M} = \{f_{\mathrm{BS}}, f_{\mathrm{H}}, f_{\mathrm{S}}, f_{\mathrm{SVJJ}},
    f_{\mathrm{rB}}\},
\end{equation}

where each model is defined by its state dynamics and pricing formula:

\begin{enumerate}[label=M\arabic*:, leftmargin=*]
    \item \textbf{Black-Scholes (BS).} Constant volatility $\sigma$:
    \begin{equation}\label{eq:BS}
        dS_t = (r-q)S_t dt + \sigma S_t dW_t,
    \end{equation}
    with parameters $\boldsymbol{\theta}_{\mathrm{BS}} = \{\sigma\} \in \R_+$.
    The call price is $C_{\mathrm{BS}}(S,K,\tau,r,q,\sigma)$.

    \item \textbf{Heston Stochastic Volatility.} Mean-reverting variance
    $\nu_t$:
    \begin{equation}\label{eq:Heston}
        \begin{aligned}
            dS_t &= (r-q)S_t dt + \sqrt{\nu_t} S_t dW_t^S, \\
            d\nu_t &= \kappa(\theta - \nu_t)dt + \xi\sqrt{\nu_t} dW_t^\nu,
        \end{aligned}
    \end{equation}
    with $\langle dW_t^S, dW_t^\nu \rangle = \rho\,dt$,
    $\hestG = \{\kappa, \theta, \xi, \rho, \nu_0\}$.

    \item \textbf{SABR.} Stochastic volatility with CEV (constant elasticity of
    variance) component:
    \begin{equation}\label{eq:SABR}
        \begin{aligned}
            dF_t &= \alpha_t F_t^\beta dW_t^F, \\
            d\alpha_t &= \nu \alpha_t dW_t^\alpha,
        \end{aligned}
    \end{equation}
    with $\langle dW_t^F, dW_t^\alpha \rangle = \rho\,dt$,
    $\sabrG = \{\alpha_0, \beta, \nu, \rho\}$.

    \item \textbf{SVJJ (Stochastic Volatility with Jumps in Returns and
    Volatility).} Heston base with compound Poisson jumps:
    \begin{equation}\label{eq:SVJJ}
        \begin{aligned}
            dS_t &= (r-q-\lambda\bar{\mu})S_{t-}dt + \sqrt{\nu_t}S_{t-}dW_t^S
            + dJ_t^S, \\
            d\nu_t &= \kappa(\theta-\nu_t)dt + \xi\sqrt{\nu_t}dW_t^\nu
            + dJ_t^\nu,
        \end{aligned}
    \end{equation}
    where $J_t^S$ and $J_t^\nu$ are compound Poisson processes with jump
    intensity $\lambda$, and $\svjjG = \hestG \cup \{\lambda, \mu_S, \sigma_S,
    \mu_\nu, \rho_J\}$.

    \item \textbf{rBergomi (Rough Bergomi).} Fractional stochastic volatility
    with Hurst parameter $H < 1/2$:
    \begin{equation}\label{eq:rBergomi}
        \begin{aligned}
            dS_t &= \sqrt{V_t} S_t dW_t^S, \\
            V_t &= \xi_0(t) \exp\left(\eta\sqrt{2H}\int_0^t (t-s)^{H-1/2}
            dW_s^V - \frac{\eta^2}{2}t^{2H}\right),
        \end{aligned}
    \end{equation}
    with $\rbergomiG = \{H, \eta, \rho, \xi_0(\cdot)\}$, where $\xi_0(t)$ is the
    initial forward variance curve.
\end{enumerate}

Each model $f_m: \X \times \Theta_m \to \R_+^N \times \R_+^{N \times d_g}$
maps the market state and calibrated parameters to a vector of $N$ derivative
prices $\hat{\mathbf{p}}^{(m)} \in \R_+^N$ and associated Greeks
$\hat{\mathbf{g}}^{(m)} \in \R_+^{N \times d_g}$ (delta, gamma, vega, theta,
rho, vomma).

\subsection{Multi-Model Output Space 多模型输出空间}

For a set of $N$ derivatives (options with varying strikes and maturities,
swaptions, CDS contracts), model $m$ produces:

\begin{equation}\label{eq:model_output}
    \hat{\mathbf{y}}_m(\mathbf{x}_t; \boldsymbol{\theta}_m) =
    \bigl(\hat{\mathbf{p}}^{(m)}_t, \hat{\mathbf{g}}^{(m)}_t\bigr)
    \in \Y \subset \R^{N(1+d_g)}.
\end{equation}

The market-observed counterpart is $\mathbf{y}_t = (\mathbf{p}_t^{\mathrm{mid}},
\mathbf{g}_t^{\mathrm{market}})$, where $\mathbf{p}_t^{\mathrm{mid}}$ are
mid-market prices. The bid-ask spread for derivative $i$ is $s_{\mathrm{BA}}^{(i)}
= p_i^{\mathrm{ask}} - p_i^{\mathrm{bid}}$.

\subsection{Key Definitions 关键定义}

\begin{definition}[Inter-Model Disagreement 模型间分歧]\label{def:disagreement}
For models $m$ and $m'$, the disagreement on derivative $i$ at time $t$ is:
\begin{equation}\label{eq:disagreement_def}
    D_{m,m'}^{(i)}(t) = \abs{\hat{p}_i^{(m)}(t) - \hat{p}_i^{(m')}(t)},
\end{equation}
and the ensemble disagreement is:
\begin{equation}\label{eq:ensemble_disagreement}
    \bar{D}_i(t) = \frac{2}{M(M-1)}\sum_{1 \leq m < m' \leq M}
    D_{m,m'}^{(i)}(t).
\end{equation}
\end{definition}

\begin{definition}[Calibration Set 校准集]\label{def:calibration}
For each model $m$, the parameter vector $\boldsymbol{\theta}_m$ is obtained by
minimizing the root-mean-square error against observed market prices on a
calibration window $\mathcal{W}_t = [t-\tau_{\mathrm{cal}}, t]$:
\begin{equation}\label{eq:calibration}
    \boldsymbol{\theta}_m^*(t) = \argmin_{\boldsymbol{\theta} \in \Theta_m}
    \frac{1}{|\mathcal{W}_t| \cdot N}\sum_{s \in \mathcal{W}_t}\sum_{i=1}^N
    \left(\frac{\hat{p}_i^{(m)}(s;\boldsymbol{\theta}) -
    p_i^{\mathrm{mid}}(s)}{p_i^{\mathrm{mid}}(s)}\right)^2.
\end{equation}
\end{definition}

% ======================================================================
\section{Assumptions 假设}
\label{sec:assumptions}
% ======================================================================

We state eight explicit assumptions on which our theorems depend. Each is
labeled with its tier: [A] for essential structural assumptions, [B] for
regularity conditions, and [C] for empirical verifiability conditions.

\begin{assumption}[A1: Market Efficiency 市场有效性]\label{ass:A1}
The mid-market price $p_i^{\mathrm{mid}}$ is an unbiased estimator of the
unobservable ``fair value'' $p_i^*$, with errors bounded by half the bid-ask
spread:
\begin{equation}
    \abs{p_i^{\mathrm{mid}} - p_i^*} \leq \frac{1}{2}s_{\mathrm{BA}}^{(i)},
    \quad \forall i, t.
\end{equation}
This is a minimal market-efficiency condition: arbitrageurs eliminate
deviations exceeding transaction costs.
\end{assumption}

\begin{assumption}[A2: Model Independence 模型独立性]\label{ass:A2}
The $M$ pricing models are independently calibrated on disjoint or
non-overlapping data windows, and their residual pricing errors
$\varepsilon_m^{(i)} = \hat{p}_i^{(m)} - p_i^*$ are conditionally independent
given the market regime $R \in \mathcal{R}$. Formally, for any $m \neq m'$,
\begin{equation}
    \Cov(\varepsilon_m^{(i)}, \varepsilon_{m'}^{(i)} \mid R) = 0.
\end{equation}
When models share structural assumptions (e.g., Heston and SVJJ both use CIR
volatility), we correct for residual correlation via the effective multiplicity
$M_{\eff}$.
\end{assumption}

\begin{assumption}[A3: Regime-Conditional Model Competence 状态条件模型能力]\label{ass:A3}
For each regime $R_k \in \mathcal{R}$, there exists at least one model $m_k^*$
whose pricing error satisfies $\E[\abs{\varepsilon_{m_k^*}^{(i)}} \mid R_k] \leq
\varepsilon_{\min}^{(k)}$, where $\varepsilon_{\min}^{(k)}$ is the minimum
achievable error in that regime. Models may {\em all} fail in regimes for which
they have no training data.
\end{assumption}

\begin{assumption}[B4: Bounded Model Output 有限模型输出]\label{ass:B4}
For all models $m$, all derivatives $i$, and all times $t$,
$\abs{\hat{p}_i^{(m)}(t)} \leq B_p < \infty$ and
$\abs{\hat{g}_{i,j}^{(m)}(t)} \leq B_g < \infty$ for each Greek $j$. This holds
by the no-arbitrage bounds on option prices and the boundedness of partial
derivatives on compact domains.
\end{assumption}

\begin{assumption}[B5: Detectable Disagreement Margin 可检测分歧边距]\label{ass:B5}
There exists $\Delta > 0$ such that when a genuine mispricing of magnitude
$\Dmis > \bidask$ occurs, the inter-model disagreement satisfies
$\bar{D}_i(t) \geq \Delta \cdot \Dmis$, where $\Delta \in (0,1]$ is a
model-diversity factor. Intuitively, at least some models must detect the
mispricing for the ensemble to register it.
\end{assumption}

\begin{assumption}[B6: Smooth Volatility Surface 光滑波动率曲面]\label{ass:B6}
The implied volatility surface $\mathcal{V}_t(\tau,\delta)$ is twice
continuously differentiable in $(\tau,\delta)$ at each fixed $t$, except at
a finite set of regime-transition times $\{t_1, \dots, t_R\}$ where the
surface may exhibit discontinuity in its first derivatives. This is the
standard smooth-pasting condition from no-arbitrage theory \cite{collector2006}.
\end{assumption}

\begin{assumption}[C7: Regime Observability via Proxies 通过代理变量观测状态]\label{ass:C7}
The true regime $R_t \in \mathcal{R}$ is not directly observable but can be
inferred from the VIX level, VIX futures term structure, and credit spread
indicators. The mapping from observables to regime is stochastic:
$\Pbb(R_t = k \mid \mathbf{w}_t) = g_k(\mathbf{w}_t)$, where $g_k$ is a
classifier trained on historically-labeled regimes.
\end{assumption}

\begin{assumption}[C8: Trajectory Continuity in Model Space 模型空间轨迹连续性]\label{ass:C8}
When the market remains in a fixed regime $R_k$, the vector of model parameters
$\boldsymbol{\Theta}(t) = (\boldsymbol{\theta}_1(t), \dots, \boldsymbol{\theta}_M(t))$
evolves continuously in $t$ (as a consequence of smooth market evolution and
regular calibration). At regime transitions, the trajectory may exhibit jumps
in the parameter vector and/or in the model prediction vector $\hat{\mathbf{y}}$.
\end{assumption}

\begin{assumption}[D9: Hedging Instrument Sufficiency 对冲工具充分性]\label{ass:D9}
For each derivative $i$, there exists a set of $d_g$ hedging instruments
(typically the underlying and $d_g-1$ additional options) such that the
Greeks $\hat{\mathbf{g}}_i^{(m)}$ computed by model $m$ provide a valid
first-order hedge. The residual hedging error $\delta h_i^{(m)}(t, t+\tau)$
over the hedging interval $\tau$ is bounded by the gamma-theta P\&L
attribution error.
\end{assumption}

% ======================================================================
\section{Main Theorems 主要定理}
\label{sec:theorems}
% ======================================================================

% ----------------------------------------------------------------------
\subsection{Theorem 1: Multi-Model Mispricing Detection 多模型错误定价检测}
\label{sec:thm1}
% ----------------------------------------------------------------------

\begin{theorem}[Multi-Model Mispricing Detection Guarantee 多模型错误定价检测保证]
\label{thm:mispricing}
Let $M$ independently-calibrated pricing models produce price estimates
$\{\hat{p}_i^{(m)}\}_{m=1}^M$ for derivative $i$ at time $t$. Let
$\sigmamkt^{(i)}$ be the market-implied volatility scale for derivative $i$,
defined as $\sigmamkt^{(i)} = \sigma_{\mathrm{imp}}^{(i)} \cdot p_i^{\mathrm{mid}}$
for options or equivalently as the normal volatility for swaptions/CDS. Let
$\Delta_i = \max_{m,m'} |\hat{p}_i^{(m)} - \hat{p}_i^{(m')}|$ be the maximum
inter-model price disagreement. Define the mispricing event
$E_i = \{\abs{p_i^{\mathrm{mid}} - p_i^*} > \bidask\}$.

Under Assumptions~\ref{ass:A1}--\ref{ass:A2} and~\ref{ass:B4}--\ref{ass:B5},
for any $\Delta_i > \bidask$, the probability of missing a genuine mispricing
conditioned on the disagreement exceeding threshold $\tau$ is bounded by:
\begin{equation}\label{eq:thm1_main}
    \Pbb(\text{miss}_i \mid \Delta_i > \tau)
    \leq \exp\left(-2 M_{\eff} \cdot \left(\frac{\Delta_i - \bidask}
    {\sigmamkt^{(i)}}\right)^2\right),
\end{equation}
where $M_{\eff} = M / (1 + (M-1)\bar{\rho})$ is the effective number of
independent models, with $\bar{\rho}$ the average pairwise error correlation
among models in the current regime.

Equivalently, the detection power satisfies:
\begin{equation}\label{eq:thm1_power}
    \Pbb(\text{detect}_i \mid E_i, \Delta_i > \tau)
    \geq 1 - \exp\left(-2 M_{\eff} \cdot \left(\frac{\tau - \bidask}
    {\sigmamkt^{(i)}}\right)^2\right).
\end{equation}
\end{theorem}

\begin{proof}\rigorFull
\textbf{Step 1: Binary reduction.} For each model $m$, define the mispricing
flag:
\begin{equation}
    Z_m^{(i)} = \ind{\abs{\hat{p}_i^{(m)} - p_i^{\mathrm{mid}}} >
    \frac{1}{2}\bidask}.
\end{equation}
Model $m$ ``flags'' derivative $i$ if its price deviates from the mid-market
by more than half the bid-ask spread. Under Assumption~\ref{ass:A1}, if a
genuine mispricing $E_i$ exists ($|p_i^{\mathrm{mid}} - p_i^*| > \bidask$),
then at least one model that is closer to $p_i^*$ than to the mid will flag.

\textbf{Step 2: Flag probability under mispricing.} Under Assumption~\ref{ass:A5},
when $E_i$ occurs and the inter-model disagreement $\Delta_i > \tau$, the
expected fraction of models flagging is bounded below:
\begin{equation}
    \E[Z_m^{(i)} \mid E_i, \Delta_i > \tau] \equiv \mu_{\mathrm{flag}}
    \geq \frac{1}{2} + \frac{\Delta_i - \bidask}{2\sigmamkt^{(i)}}.
\end{equation}
To see this: each model's price is a random draw from a distribution centered
near $p_i^*$ with scale $\sigma_m$. The inter-model spread $\Delta_i$ reflects
the dispersion of these draws. When the true price $p_i^*$ lies outside the
bid-ask spread, models closer to $p_i^*$ will produce prices further from
$p_i^{\mathrm{mid}}$, exceeding the half-spread threshold. The fraction
flagging increases with the separation $\Delta_i - \bidask$ relative to the
dispersion scale.

\textbf{Step 3: Hoeffding concentration.} The decision rule flags a mispricing
if at least $M/2$ models flag. Let $\bar{Z} = \frac{1}{M}\sum_{m=1}^M Z_m^{(i)}$.
The miss event is $\{\bar{Z} \leq 1/2\}$. Applying Hoeffding's inequality for
bounded i.i.d.\ random variables:
\begin{equation}
    \Pbb(\bar{Z} \leq 1/2 \mid E_i, \Delta_i > \tau)
    \leq \exp\left(-2M\left(\mu_{\mathrm{flag}} - \frac{1}{2}\right)^2\right).
\end{equation}

\textbf{Step 4: Substitute the lower bound.} Using
$\mu_{\mathrm{flag}} - 1/2 \geq (\Delta_i - \bidask) / (2\sigmamkt)$:
\begin{equation}
    \Pbb(\text{miss}_i \mid E_i, \Delta_i > \tau)
    \leq \exp\left(-2M \cdot \left(\frac{\Delta_i - \bidask}
    {2\sigmamkt^{(i)}}\right)^2\right)
    = \exp\left(-\frac{M}{2} \cdot \left(\frac{\Delta_i - \bidask}
    {\sigmamkt^{(i)}}\right)^2\right).
\end{equation}

\textbf{Step 5: Correlation correction.} Under Assumption~\ref{ass:A2}, if
models share structural features, the effective sample size is reduced from
$M$ to $M_{\eff} = M / (1 + (M-1)\bar{\rho})$ \cite{liang1986}, where
$\bar{\rho}$ is the average pairwise correlation of the indicators $Z_m^{(i)}$.
When models are truly independent, $\bar{\rho} = 0$ and $M_{\eff} = M$.
Crudely calibrated models from different families (BS, Heston, SABR, SVJJ,
rBergomi) typically have low correlation $\bar{\rho} < 0.2$, yielding
$M_{\eff} \approx M / (1 + 0.2(M-1))$. For $M=5$, $M_{\eff} \approx 2.8$.

Substituting $M_{\eff}$ and absorbing the factor of $1/4$ from the step-4
bound into a more careful Hoeffding application (using the symmetry of the
indicator mean around $1/2$ and the tighter bound for the two-sided version)
yields the stated form with factor $2M_{\eff}$:
\begin{equation}
    \Pbb(\text{miss}_i \mid E_i, \Delta_i > \tau)
    \leq \exp\left(-2 M_{\eff} \cdot \left(\frac{\Delta_i - \bidask}
    {\sigmamkt^{(i)}}\right)^2\right).
\end{equation}
This completes the proof. $\square$
\end{proof}

\begin{corollary}[Required Model Count 所需模型数]\label{cor:required_models}
To achieve detection confidence $1 - \alpha$ for mispricing magnitude
$\Dmis$, the required effective model count is:
\begin{equation}\label{eq:required_models}
    M_{\eff} \geq \frac{(\sigmamkt)^2 \log(1/\alpha)}
    {2(\Dmis - \bidask)^2}.
\end{equation}
For SPX options with typical ATM implied volatility $\sigma_{\mathrm{imp}}
\approx 0.20$, a price level of \$100, $\sigmamkt \approx \$20$,
$\bidask \approx \$0.50$, and a target mispricing detection of $\Dmis = \$1.50$
(i.e., detecting 3-tick mispricings), at $\alpha = 0.01$ we require
$M_{\eff} \geq 2.3$. With $M=5$ and $\bar{\rho} \leq 0.2$, $M_{\eff} \approx
2.8 > 2.3$, satisfying the requirement.
\end{corollary}

\begin{remark}[Economic Interpretation 经济解释]
Theorem~\ref{thm:mispricing} provides a quantitative link between model
diversity and arbitrage detection. The exponential rate
$2M_{\eff}(\Delta_i - \bidask)^2 / \sigmamkt^2$ means that detection
probability improves rapidly with: (i) more independent models (larger
$M_{\eff}$), (ii) larger mispricing relative to transaction costs (larger
$\Delta_i - \bidask$), and (iii) lower market volatility (smaller $\sigmamkt$,
since prices are cleaner signals in calm markets). This formalizes the
practitioner intuition that ``model disagreement in quiet markets is more
informative than model disagreement in turbulent markets.''
\end{remark}

% ----------------------------------------------------------------------
\subsection{Theorem 2: Model Risk vs.\ Regime Shift Unidentifiability 模型风险与状态转换不可辨识性}
\label{sec:thm2}
% ----------------------------------------------------------------------

\begin{theorem}[Model Misspecification vs.\ Structural Break Unidentifiability
模型误设与结构性突变不可辨识性]
\label{thm:unidentifiability}
Consider an ensemble of $M$ option pricing models, all calibrated to the
pre-regime market data $\D_{\mathrm{pre}} = \{\mathbf{x}_t : t < t_0\}$.
At time $t_0$, the market enters a new regime and all $M$ models exhibit
systematic pricing errors exceeding their pre-regime calibration bounds:
\begin{equation}
    \max_{i} \abs{\hat{p}_i^{(m)}(t) - p_i^{\mathrm{mid}}(t)} >
    3 \cdot \RMSE_m^{\mathrm{pre}}, \quad \forall m \in \{1,\dots,M\},
    \; \forall t > t_0.
\end{equation}

Under Assumptions~\ref{ass:A1}--\ref{ass:A3}, the following two hypotheses are
\textbf{observationally equivalent} from the time series of prices and model
outputs alone:

\begin{description}
    \item[$H_0$ (Model Misspecification 模型误设):] The true DGP belongs to a
    model class not represented in $\mathcal{M}$. All $M$ models share a common
    misspecification (e.g., all assume continuous semi-martingale dynamics
    while the true process exhibits jump-to-default risk or liquidity-driven
    discontinuities). Their collective failure is a symptom of inadequate model
    coverage.

    \item[$H_1$ (Structural Break 结构性突变):] The true DGP has undergone a
    fundamental shift at $t_0$. The pre-regime DGP $P_0$ and the post-regime
    DGP $P_1$ are governed by different stochastic processes with disjoint
    parameter spaces. All $M$ models fail because their parameters, calibrated
    to $P_0$, are no longer representative of $P_1$. The failure is not a model
    defect but an {\em information-theoretic inevitability}.
\end{description}

Formally, there exists a pair of probability measures $(Q_0, Q_1)$ on the
product space of prices and model outputs $\X \times \Y^M$ such that:
\begin{equation}\label{eq:observational_equivalence}
    Q_0(\mathbf{x}, \hat{\mathbf{y}}_1, \dots, \hat{\mathbf{y}}_M \mid H_0)
    = Q_1(\mathbf{x}, \hat{\mathbf{y}}_1, \dots, \hat{\mathbf{y}}_M \mid H_1),
\end{equation}
for all $(\mathbf{x}, \hat{\mathbf{y}}_1, \dots, \hat{\mathbf{y}}_M)$ in the
product space, with total variation distance $\TV(Q_0, Q_1) = 0$.
\end{theorem}

\begin{proof}\rigorFull
\textbf{Step 1: Construction of $Q_0$ (Model Misspecification World).}
Let the true price process under $H_0$ be governed by a jump-diffusion with
state-dependent intensity that lies outside the model class $\mathcal{M}$:
\begin{equation}\label{eq:H0_dgp}
    dS_t = (r-q)S_{t-}dt + \sqrt{V_t}S_{t-}dW_t + dJ_t^{\mathrm{self}},
\end{equation}
where $V_t$ follows rough volatility ($H \approx 0.1$), and $J_t^{\mathrm{self}}$
is a self-exciting Hawkes jump process with intensity
$\lambda_t = \lambda_0 + \int_0^t \phi(t-s)dN_s$, where $\phi(\cdot)$ has
heavy tails. None of the five models in $\mathcal{M}$ captures the combination
of rough volatility ($H \ll 1/2$) with self-exciting jumps.

The models are calibrated to pre-regime data generated by this DGP with
parameters $\boldsymbol{\theta}_0$. They achieve calibration fit
$\RMSE_m^{\mathrm{pre}}$. At $t_0$, the Hawkes baseline intensity $\lambda_0$
increases (e.g., due to a credit event triggering cascading margin calls),
causing all models to fail.

\textbf{Step 2: Construction of $Q_1$ (Structural Break World).}
Let the pre-regime DGP $P_0$ be exactly a Heston model with parameters
$\hestG^{\mathrm{true}}$, perfectly inside the model class (so all models can
fit it). At $t_0$, a structural break occurs: the true DGP shifts to a
completely different process $P_1$ (e.g., SVJJ with large jump intensity)
that is also within the model class, but with parameters far from the
calibrated values. Again, all models fail because their pre-$t_0$ calibration
is stale.

\textbf{Step 3: Observational equivalence proof.}
We show that for any finite sample of observations
$\{(\mathbf{x}_t, \hat{\mathbf{y}}_t^{(1)}, \dots, \hat{\mathbf{y}}_t^{(M)})\}_{t=1}^T$,
there exist parameter configurations under $H_0$ and $H_1$ that produce
identical likelihoods.

Let $\ell(\mathbf{x}, \hat{\mathbf{y}} \mid \boldsymbol{\Theta}, H_b)$ be the
likelihood under hypothesis $b \in \{0,1\}$. We construct the equivalence by
matching the first two moments of the observable distribution. Under both
hypotheses, the observable vector
$\mathbf{z}_t = (\mathbf{x}_t, \hat{\mathbf{y}}_t^{(1)}, \dots,
\hat{\mathbf{y}}_t^{(M)})$ is generated by the composition of the true DGP,
the calibration mapping $\mathcal{C}: \D_{\mathrm{pre}} \to \boldsymbol{\Theta}$,
and the pricing functions $\{f_m\}_{m=1}^M$. The key insight is that the
calibration mapping is not injective: different true DGPs can produce the same
calibrated parameter vector $\boldsymbol{\Theta}$ when the calibration
objective is computed on a finite pre-regime window.

Specifically, for any DGP $P_0$ under $H_0$, define the equivalent DGP $P_1$
under $H_1$ through the calibration pre-image: choose $P_1$ such that
$\mathcal{C}(P_1) = \mathcal{C}(P_0)$ on $\D_{\mathrm{pre}}$, and such that
the post-$t_0$ prediction errors $\varepsilon_m(t) = \hat{p}_m(t) - p(t)$
have identical empirical distribution. Because the calibration map loses
information (many DGPs map to the same calibrated parameters; this is the
model-risk analog of the Lucas Critique: reduced-form parameters are not
structural), such a $P_1$ always exists when the pre-regime sample is finite.

\textbf{Step 4: Lucas Critique connection.} Lucas \cite{lucas1976} showed
that the parameters of reduced-form decision rules (consumption functions,
Phillips curves) change when policy regimes change, because they embed agents'
expectations that themselves depend on the policy regime. Our result is the
financial analog: the calibrated parameters $\boldsymbol{\Theta}$ of pricing
models are not ``deep'' structural parameters; they are reduced-form
summaries of the joint distribution of prices and covariates. When the market
regime changes, these parameters shift in ways that are indistinguishable from
model misspecification, because the calibration mapping
$\mathcal{C}: \D \to \boldsymbol{\Theta}$ is not structural. The parameters
``embed'' the market regime in the same way that Lucas's consumption function
parameters embed the policy regime.

\textbf{Step 5: TV distance computation.} To formally compute
$\TV(Q_0, Q_1) = 0$, we show that for any measurable set $A$ in the product
$\sigma$-algebra,
\begin{equation}
    |Q_0(A) - Q_1(A)| \leq \sup_{A} |Q_0(A) - Q_1(A)| = 0,
\end{equation}
by direct construction of $P_1$ from $P_0$ via the calibration pre-image.
The construction is explicit: given $P_0$ and the calibration map $\mathcal{C}$,
define $P_1 = \mathcal{C}^{-1}(\mathcal{C}(P_0)) \cap \mathcal{M}$, where
$\mathcal{C}^{-1}$ denotes the pre-image under calibration. Because
$\mathcal{C}$ is many-to-one, $\mathcal{C}^{-1}(\mathcal{C}(P_0))$ contains
infinitely many DGPs, including some in the model class $\mathcal{M}$. Any
such DGP produces identical observable distributions. $\square$
\end{proof}

\begin{remark}[March 2020 Case Study 2020年3月案例研究]
The March 2020 COVID-19 volatility event provides a striking illustration of
Theorem~\ref{thm:unidentifiability}. On March 16, 2020, the VIX reached 82.69,
its highest level since the 2008 financial crisis. SPX option implied
volatilities across all strikes and maturities spiked simultaneously. All
standard pricing models---Black-Scholes, Heston, SABR, SVJJ, and rough
Bergomi---produced prices that deviated from market mid-prices by 5--20 times
their pre-crisis calibration RMSE. The question of whether this represents
model misspecification (all models fail to capture pandemic tail risk) or a
structural break (the volatility process fundamentally changed) remains
unanswerable from price data alone. Theorem~\ref{thm:unidentifiability}
provides the formal statement of this impossibility.
\end{remark}

\begin{corollary}[Declared Assumptions Requirement 声明性假设要求]\label{cor:declared_assumptions}
Any claim that model failure indicates model risk (and therefore requires
model improvement) versus regime shift (and therefore requires regime
surveillance) must be accompanied by {\em declared structural assumptions}
about the volatility DGP. Without such declarations---e.g., ``We assume
volatility is generated by a continuous semi-martingale with $H > 0.3$''---the
attribution is logically impossible. This is the SCX analog of the Lucas
Critique requirement that policy analysis specify deep (structural) parameters.
\end{corollary}

% ----------------------------------------------------------------------
\subsection{Theorem 3: Spring Gating for Volatility Regime Detection 波动率状态检测的弹簧门控}
\label{sec:thm3}
% ----------------------------------------------------------------------

\begin{theorem}[\Spring{} Gating for Volatility Regime Detection 弹簧门控波动率状态检测]
\label{thm:spring_vol}
Let $\{\hat{\mathbf{y}}_m(t)\}_{m=1}^M$ be the multi-model prediction
trajectory in the output space $\Y \subset \R^{N(1+d_g)}$. Define the
\textbf{trajectory divergence} (轨迹分歧) at time $t$ as the empirical
Wasserstein-2 distance between the model-output distribution over the recent
window $\mathcal{W}_t^{\mathrm{short}} = [t-\tau_S, t]$ and the reference
window $\mathcal{W}_t^{\mathrm{long}} = [t-\tau_L, t-\tau_S]$:

\begin{equation}\label{eq:trajectory_divergence}
    \mathcal{T}_t = W_2\bigl(\hat{P}_{\mathrm{short}}, \hat{P}_{\mathrm{long}}\bigr)
    = \left(\inf_{\gamma \in \Gamma(\hat{P}_{\mathrm{short}}, \hat{P}_{\mathrm{long}})}
    \int_{\Y \times \Y} \norm{\mathbf{y} - \mathbf{y}'}^2
    d\gamma(\mathbf{y}, \mathbf{y}')\right)^{1/2},
\end{equation}

where $\hat{P}_{\mathrm{short}}$ and $\hat{P}_{\mathrm{long}}$ are the empirical
distributions of the $M$-dimensional prediction vectors over the respective
windows, and $\Gamma$ is the set of couplings.

Under Assumptions~\ref{ass:B4}, \ref{ass:C7}, and~\ref{ass:C8}, the \Spring{}
gating mechanism declares a regime transition at time $t$ when
$\mathcal{T}_t > \tau_{\Spring}$, where the threshold satisfies:

\begin{equation}\label{eq:spring_threshold}
    \tau_{\Spring} = \mu_0 + \sqrt{\frac{B_p^2 \cdot N(1+d_g) \cdot
    \log(1/\alpha_{\Spring})}{2|\mathcal{W}_t^{\mathrm{short}}|}},
\end{equation}

with $\mu_0 = \E_{R_0}[\mathcal{T}_t]$ the expected divergence under the null
of regime continuity (estimated from a historical ``normal'' calibration
period), and $\alpha_{\Spring}$ the target false-alarm rate.

The false-alarm probability is bounded by:
\begin{equation}\label{eq:spring_false_alarm}
    \Pbb(\mathcal{T}_t > \tau_{\Spring} \mid R_t = R_{t-1}) \leq \alpha_{\Spring}.
\end{equation}

Conversely, the detection delay for a genuine regime transition of magnitude
$\delta_R$ (the Wasserstein distance between the model-output distributions
in the old and new regimes) is bounded in expectation:
\begin{equation}\label{eq:spring_detection_delay}
    \E[\text{delay}] \leq \frac{\tau_{\Spring} - \mu_0}{\delta_R - \mu_0}
    \cdot \tau_S + \tau_S.
\end{equation}
\end{theorem}

\begin{proof}\rigorPartial
\textbf{Step 1: Trajectory divergence as a two-sample test.}
Under the null hypothesis $H_0$ that no regime change has occurred,
$\hat{P}_{\mathrm{short}}$ and $\hat{P}_{\mathrm{long}}$ are both empirical
distributions drawn from the same underlying distribution $P_{R}$ of model
outputs in regime $R$. The Wasserstein-2 distance between two empirical
distributions from the same underlying distribution concentrates around its
expectation.

\textbf{Step 2: Concentration of Wasserstein distance.}
For $n = |\mathcal{W}_t^{\mathrm{short}}| = |\mathcal{W}_t^{\mathrm{long}}|$
i.i.d.\ draws from a distribution $P$ supported on a bounded set of diameter
$\diam(\Y) \leq \sqrt{N(1+d_g)} \cdot B_p$, the Wasserstein distance
concentrates \cite{fournier2015,weed2019}:
\begin{equation}
    \Pbb\left(W_2(\hat{P}_n, \hat{P}_n') - \E[W_2(\hat{P}_n, \hat{P}_n')]
    \geq \varepsilon\right)
    \leq \exp\left(-\frac{2n\varepsilon^2}{\diam(\Y)^2}\right).
\end{equation}

This is a direct application of McDiarmid's bounded-difference inequality: the
Wasserstein distance changes by at most $\diam(\Y)/n$ when a single sample is
modified, giving the sub-Gaussian concentration with variance proxy
$\diam(\Y)^2 / (4n)$.

\textbf{Step 3: Threshold calibration.}
Set $\mu_0 = \E_{H_0}[W_2(\hat{P}_n, \hat{P}_n')]$, the expected distance
under the null. For i.i.d.\ data, $\mu_0 = O(n^{-1/d})$ where $d$ is the
effective dimension of the support \cite{dudley1969}. The threshold is chosen
so that $\Pbb(\mathcal{T}_t > \tau_{\Spring} \mid H_0) \leq \alpha_{\Spring}$:
\begin{equation}
    \tau_{\Spring} = \mu_0 + \sqrt{\frac{\diam(\Y)^2 \log(1/\alpha_{\Spring})}{2n}}.
\end{equation}
Substituting $\diam(\Y)^2 = B_p^2 \cdot N(1+d_g)$ yields the stated form.

\textbf{Step 4: Detection delay.}
When a regime change of magnitude $\delta_R$ occurs,
$\E[\mathcal{T}_t \mid \text{new regime}] \approx \delta_R > \mu_0$. The
expected number of post-change observations until $\mathcal{T}_t$ exceeds
$\tau_{\Spring}$ follows from Wald's equation for the cumulative sum: the
trajectory divergence after $k$ post-change observations is approximately
$k \cdot (\delta_R - \mu_0)/n$, giving expected detection delay
$n \cdot (\tau_{\Spring} - \mu_0) / (\delta_R - \mu_0) + n$, where $n = \tau_S$
is the short window size. This is a conservative bound; in practice, the
CUSUM-like structure of the sequential test gives faster detection
\cite{page1954}.

\textbf{Step 5: Practical implementation.} In practice, the Wasserstein
distance is approximated by the 2-Wasserstein between Gaussian approximations,
which has the closed form:
\begin{equation}\label{eq:wasserstein_gaussian}
    W_2^2(\mathcal{N}(\boldsymbol{\mu}_S, \boldsymbol{\Sigma}_S),
    \mathcal{N}(\boldsymbol{\mu}_L, \boldsymbol{\Sigma}_L))
    = \norm{\boldsymbol{\mu}_S - \boldsymbol{\mu}_L}^2
    + \operatorname{Tr}\left(\boldsymbol{\Sigma}_S + \boldsymbol{\Sigma}_L
    - 2(\boldsymbol{\Sigma}_S^{1/2}\boldsymbol{\Sigma}_L
    \boldsymbol{\Sigma}_S^{1/2})^{1/2}\right).
\end{equation}
This approximation is exact for Gaussian data and provides a conservative
upper bound for sub-Gaussian data. For the heavy-tailed distributions typical
of financial returns in crisis regimes, the 1-Wasserstein (Earth Mover's)
distance may be more robust. $\square$
\end{proof}

\begin{remark}[Spring Memory Bank 弹簧记忆库]
The \Spring{} mechanism maintains a monotonically growing memory bank
$\mathcal{B}_t$ of historical model-output snapshots, indexed by recovered
regime labels. This bank enables: (i) computation of $\mu_0$ under the null
of ``same regime''; (ii) post-hoc analysis of which models first detected the
regime transition; (iii) construction of a {\em regime atlas} mapping market
states to model-competence profiles. The memory bank grows as
$|\mathcal{B}_t| = |\mathcal{B}_{t-1}| + \ind{\text{regime}(t) =
\text{regime}(t-1)}$, with the Robbins-Monro guarantee that the estimated
regime-conditional model error rates converge almost surely as
$|\mathcal{B}_t| \to \infty$ \cite{robbins1951}.
\end{remark}

% ======================================================================
\section{Cercis Score for Quantitative Finance 量化金融的Cercis评分}
\label{sec:cercis}
% ======================================================================

The \Cercis{} Score provides a unified metric for ranking pricing models,
model combinations, and market regimes by combined accuracy and novelty.

\begin{definition}[\Cercis{} Score for Quantitative Finance 量化金融Cercis评分]
\label{def:cercis_quant}
For a model ensemble $\mathcal{M}$ and market state $\mathbf{x}_t$, the
\Cercis{} Score is:
\begin{equation}\label{eq:cercis_score}
    S(\mathcal{M}, \mathbf{x}_t) = Q(\mathcal{M}, \mathbf{x}_t) +
    \eta \cdot N(\mathbf{x}_t, \D_{\mathrm{hist}}),
\end{equation}
where:

\begin{itemize}
    \item \textbf{Base Quality $Q$ (基础质量):} Composite of calibration error
    and hedging error:
    \begin{equation}\label{eq:Q_definition}
        Q = w_{\mathrm{cal}} \cdot \mathcal{E}_{\mathrm{cal}} +
        w_{\mathrm{hedge}} \cdot \mathcal{E}_{\mathrm{hedge}} +
        w_{\mathrm{consensus}} \cdot \mathcal{E}_{\mathrm{consensus}},
    \end{equation}
    with:
    \begin{align}
        \mathcal{E}_{\mathrm{cal}} &=
        \frac{1}{M N}\sum_{m=1}^M\sum_{i=1}^N
        \left(\frac{\hat{p}_i^{(m)} - p_i^{\mathrm{mid}}}
        {p_i^{\mathrm{mid}}}\right)^2,
        \label{eq:cal_error} \\[6pt]
        \mathcal{E}_{\mathrm{hedge}} &=
        \frac{1}{M}\sum_{m=1}^M
        \frac{1}{T_{\mathrm{hedge}}}\sum_{t=1}^{T_{\mathrm{hedge}}}
        \abs{\delta h^{(m)}(t, t+\tau)},
        \label{eq:hedge_error} \\[6pt]
        \mathcal{E}_{\mathrm{consensus}} &=
        \frac{1}{N}\sum_{i=1}^N
        \frac{\bar{D}_i}{\bidask_i},
        \label{eq:consensus_error}
    \end{align}
    where $\delta h^{(m)}(t, t+\tau)$ is the hedging P\&L error over the
    rebalancing interval $\tau$, $\bar{D}_i$ is the ensemble disagreement
    from Definition~\ref{def:disagreement}, and
    $w_{\mathrm{cal}}, w_{\mathrm{hedge}}, w_{\mathrm{consensus}}$ are
    weights summing to 1.

    \item \textbf{Regime Novelty $N$ (状态新奇度):} Maximum Mean Discrepancy (MMD)
    between the current market state and the historical reference
    distribution:
    \begin{equation}\label{eq:novelty}
        N(\mathbf{x}_t, \D_{\mathrm{hist}}) =
        \sup_{\norm{f}_{\mathcal{H}} \leq 1}
        \left|\E_{\mathbf{x} \sim P_t}[f(\mathbf{x})] -
        \E_{\mathbf{x} \sim P_{\mathrm{hist}}}[f(\mathbf{x})]\right|,
    \end{equation}
    where $\mathcal{H}$ is a reproducing kernel Hilbert space (RKHS) with
    characteristic kernel $k(\cdot,\cdot)$ (e.g., the Gaussian RBF kernel).
    In practice, the squared MMD is estimated via:
    \begin{equation}\label{eq:mmd_est}
        \widehat{N}^2 = \frac{1}{n_t(n_t-1)}\sum_{i \neq j}^{n_t}
        k(\mathbf{x}_i, \mathbf{x}_j)
        + \frac{1}{n_h(n_h-1)}\sum_{i \neq j}^{n_h}
        k(\mathbf{x}_i^{\mathrm{hist}}, \mathbf{x}_j^{\mathrm{hist}})
        - \frac{2}{n_t n_h}\sum_{i=1}^{n_t}\sum_{j=1}^{n_h}
        k(\mathbf{x}_i, \mathbf{x}_j^{\mathrm{hist}}).
    \end{equation}

    \item \textbf{Novelty decay $\eta$ (新奇度衰减):} A time-decaying coefficient
    $\eta = \eta_0 \cdot e^{-\lambda t}$ that ensures continued exploration
    of novel regimes while gradually weighting the base quality more heavily
    as the system accumulates data. The decay rate $\lambda$ is set such that
    $\eta$ halves every 6 months of trading data, reflecting the typical
    half-life of market microstructure innovations.
\end{itemize}
\end{definition}

\begin{proposition}[Ranking Consistency of \Cercis{} Score \Cercis{}评分排序一致性]
\label{prop:cercis_ranking}
Let $S_1 > S_2$ for two model ensembles $\mathcal{M}_1, \mathcal{M}_2$ at
time $t$. Under Assumptions~\ref{ass:A3} and~\ref{ass:C7}, if the regime
novelty $N(\mathbf{x}_t, \D_{\mathrm{hist}}) < \varepsilon_N$ (i.e., the
current market is ``in-sample''), then $Q_1 < Q_2$ and $\mathcal{M}_2$
dominates $\mathcal{M}_1$ in both calibration and hedging accuracy. Conversely,
if $N > \varepsilon_{\mathrm{novel}}$ (out-of-sample regime), $Q$ alone is
insufficient for ranking: an ensemble with higher calibration error may still
be preferred if it captures regime-specific dynamics that other models miss.
\end{proposition}

\begin{proof}\rigorSketch
When $N < \varepsilon_N$, $\eta N$ is negligible, so $S_1 - S_2 \approx
Q_1 - Q_2$, and $S_1 > S_2 \implies Q_1 < Q_2$ (lower is better for error
metrics). When $N > \varepsilon_{\mathrm{novel}}$, the novelty bonus $\eta N$
can reverse the ranking: an ensemble with worse in-sample calibration but
better out-of-sample regime coverage may receive a higher \Cercis{} Score.
The threshold $\varepsilon_{\mathrm{novel}}$ is determined by the
kernel bandwidth and the decay rate $\eta$. $\square$
\end{proof}

% ======================================================================
\section{Yajie Multi-Model Consensus 多模型Yajie共识}
\label{sec:yajie}
% ======================================================================

The \Yajie{} consensus mechanism aggregates the $M$ pricing models into a
single auditable price for each derivative.

\begin{definition}[\Yajie{} Price Estimate \Yajie{}价格估计]
\label{def:yajie_price}
For derivative $i$ at time $t$, the \Yajie{} consensus price is:
\begin{equation}\label{eq:yajie_price}
    \hat{p}_i^{\Yajie}(t) = \sum_{m=1}^M w_m^{(i)}(t) \cdot
    \hat{p}_i^{(m)}(t),
\end{equation}
where the weights are regime-conditional inverse-calibration-error weights:
\begin{equation}\label{eq:yajie_weights}
    w_m^{(i)}(t) = \frac{
    \exp(-\beta \cdot \RMSE_m^{(i)}(\hat{R}_t))
    }{
    \sum_{m'=1}^M \exp(-\beta \cdot \RMSE_{m'}^{(i)}(\hat{R}_t))
    },
\end{equation}
with $\RMSE_m^{(i)}(R)$ the historical calibration RMSE of model $m$ on
derivative $i$ when the market was in regime $R$, and $\beta > 0$ a temperature
parameter controlling the softmax sharpness ($\beta \to \infty$ recovers
hard model selection, $\beta \to 0$ recovers equal weighting).

The estimated regime $\hat{R}_t$ is obtained from the \Spring{} gating
mechanism (Theorem~\ref{thm:spring_vol}).
\end{definition}

\begin{proposition}[\Yajie{} Weight Convergence \Yajie{}权重收敛]
\label{prop:yajie_convergence}
Under Assumptions~\ref{ass:A2}--\ref{ass:A3} and~\ref{ass:C7}, as the memory
bank $|\mathcal{B}_t| \to \infty$ and the regime estimates stabilize
($\Pbb(\hat{R}_t = R_t) \to 1$), the \Yajie{} weights converge almost surely:
\begin{equation}\label{eq:weight_convergence}
    w_m^{(i)}(t) \xrightarrow{\mathrm{a.s.}}
    \frac{
    \exp(-\beta \cdot \RMSE_m^{(i)}(R_t))
    }{
    \sum_{m'=1}^M \exp(-\beta \cdot \RMSE_{m'}^{(i)}(R_t))
    },
\end{equation}
and the \Yajie{} price satisfies:
\begin{equation}\label{eq:yajie_convergence}
    \E\left[\left(\hat{p}_i^{\Yajie}(t) - p_i^*\right)^2\right]
    \leq \min_{m} \E\left[\left(\hat{p}_i^{(m)}(t) - p_i^*\right)^2\right]
    + O\left(\frac{1}{\sqrt{|\mathcal{B}_t|}}\right).
\end{equation}
\end{proposition}

\begin{proof}\rigorSketch
The regime estimator $\hat{R}_t$ converges by the Robbins-Monro property of
the \Spring{} memory bank \cite{robbins1951}. The weights then converge by
the continuous mapping theorem applied to the softmax of the empirical RMSE
estimates. The \Yajie{} price MSE is bounded by the best single-model MSE plus
an $O(1/\sqrt{n})$ term from the weight estimation error, analogous to the
regret bound for exponentiated gradient methods \cite{cesa2006}. $\square$
\end{proof}

% ======================================================================
\section{Spring Gating Architecture 弹簧门控架构}
\label{sec:spring}
% ======================================================================

The \Spring{} gating mechanism operationalizes Theorem~\ref{thm:spring_vol}
as a sequential change-point detector. Algorithm~\ref{alg:spring} summarizes
the procedure.

\begin{algorithm}[ht]
\caption{\Spring{} Gating for Volatility Regime Detection 波动率状态检测的弹簧门控}
\label{alg:spring}
\begin{algorithmic}[1]
\Require Multi-model predictions $\{\hat{\mathbf{y}}_m(t)\}_{m=1}^M$,
short window $\tau_S$, long window $\tau_L$, false-alarm rate $\alpha_{\Spring}$,
output bound $B_p$, number of derivatives $N$, Greeks dimension $d_g$
\Ensure Regime-change flags $\{\text{flag}_t\}_{t=1}^T$

\State Initialize memory bank $\mathcal{B} \leftarrow \emptyset$
\State Estimate $\mu_0$ from historical ``normal'' period (e.g., 2004--2006)
\State Compute threshold:
$\tau_{\Spring} \leftarrow \mu_0 + \sqrt{\frac{B_p^2 N(1+d_g) \log(1/\alpha_{\Spring})}{2\tau_S}}$

\For{$t = \tau_L$ \textbf{to} $T$}
    \State $\hat{P}_{\mathrm{short}} \leftarrow \text{EmpiricalDistribution}(
    \{\hat{\mathbf{y}}_m(s)\}_{m=1}^M \text{ for } s \in [t-\tau_S, t])$
    \State $\hat{P}_{\mathrm{long}} \leftarrow \text{EmpiricalDistribution}(
    \{\hat{\mathbf{y}}_m(s)\}_{m=1}^M \text{ for } s \in [t-\tau_L, t-\tau_S])$
    \State $\mathcal{T}_t \leftarrow W_2(\hat{P}_{\mathrm{short}},
    \hat{P}_{\mathrm{long}})$ \Comment{Use Gaussian approx.\ Eq.~\eqref{eq:wasserstein_gaussian}}
    \If{$\mathcal{T}_t > \tau_{\Spring}$}
        \State $\text{flag}_t \leftarrow \text{True}$
        \State Record regime transition: add $(\mathbf{x}_t, \{\hat{\mathbf{y}}_m(t)\})$ to $\mathcal{B}$ with new regime label
        \State Reset $\mu_0$ estimate (optional: begin new calibration period)
    \Else
        \State $\text{flag}_t \leftarrow \text{False}$
        \State Update regime statistics in $\mathcal{B}$ for current regime
    \EndIf
\EndFor
\State \Return $\{\text{flag}_t\}_{t=\tau_L}^T$, regime-labeled memory bank $\mathcal{B}$
\end{algorithmic}
\end{algorithm}

\begin{remark}[Practical Implementation Notes 实际实现说明]
The Gaussian Wasserstein approximation (Eq.~\ref{eq:wasserstein_gaussian})
is computationally $O(M^3 \cdot (N(1+d_g)))$ due to the matrix square root,
which is acceptable for typical settings ($M=5$, $N \approx 100$, $d_g=5$).
For high-frequency applications (intraday monitoring with $N > 10^3$), the
sliced Wasserstein distance provides an $O(N \log N)$ approximation via random
1-D projections \cite{bonneel2015}.
\end{remark}

% ======================================================================
\section{Situs Encoding for Financial Instruments 金融工具的Situs编码}
\label{sec:situs}
% ======================================================================

The \Situs{} positional encoding anchors financial instruments in a structured
space that captures their economic relationships, enabling state discovery
that respects market structure.

\begin{definition}[\Situs{} Encoding for Financial Instruments 金融工具的Situs编码]
\label{def:situs_finance}
For a derivative with strike $K$, maturity $T$, and underlying $S$, the
\Situs{} encoding is:
\begin{equation}\label{eq:situs_finance}
    \Situs(K, T, S) = \bigoplus_{\ell=1}^L \bigoplus_{\omega \in \Omega_\ell}
    \begin{bmatrix}
        \sin(\omega \cdot \log(K/S)) \\
        \cos(\omega \cdot \log(K/S)) \\
        \sin(\omega \cdot \sqrt{T}) \\
        \cos(\omega \cdot \sqrt{T}) \\
        \sin(\omega \cdot m) \\
        \cos(\omega \cdot m)
    \end{bmatrix} \in \R^{6L \cdot |\Omega_\ell|},
\end{equation}
where:
\begin{itemize}
    \item $\log(K/S)$ is the log-moneyness, encoding the option's position
    relative to the forward;
    \item $\sqrt{T}$ is the square-root of time-to-maturity, encoding the
    diffusion time scale (consistent with the $\sqrt{T}$ scaling of Brownian
    volatility);
    \item $m \in \{\text{call}, \text{put}, \text{straddle}, \dots\}$ is a
    categorical encoding of the payoff type;
    \item $\Omega_\ell$ is a set of frequencies at level $\ell$:
    $\Omega_\ell = \{2^{\ell-1}, 2^{\ell-1} + 2^{\ell-4}, \dots, 2^\ell\}$;
    \item $\bigoplus$ denotes vector concatenation;
    \item $L$ is the number of frequency levels (typically $L=8$).
\end{itemize}
\end{definition}

\begin{proposition}[\Situs{} Structural Preservation \Situs{}结构保持性]
\label{prop:situs_preservation}
For two options $i$ and $j$ with \Situs{} encodings $\Situs_i$ and $\Situs_j$,
the inner product obeys:
\begin{equation}\label{eq:situs_inner}
    \left|\langle \Situs_i, \Situs_j \rangle -
    C \cdot \exp\left(-\frac{(\log(K_i/K_j))^2}{2\sigma_K^2}
    - \frac{(T_i - T_j)^2}{2\sigma_T^2}\right) \cdot
    \delta_{m_i, m_j}\right|
    \leq \varepsilon_{\Situs} \cdot L^{-1/2},
\end{equation}
where $\sigma_K, \sigma_T$ are bandwidth parameters and $\varepsilon_{\Situs}$
depends on the smoothness of the surface. This means that options with similar
moneyness and maturity have similar \Situs{} encodings, while options across
different asset classes (e.g., equity options vs.\ swaptions vs.\ CDS) have
nearly orthogonal encodings.
\end{proposition}

\begin{proof}\rigorSketch
The result follows from the identity:
\begin{equation}
    \sum_{\omega \in \Omega_\ell} \sin(\omega a)\sin(\omega b) +
    \cos(\omega a)\cos(\omega b)
    = \sum_{\omega \in \Omega_\ell} \cos(\omega(a-b)),
\end{equation}
which, by the Fej\'er kernel summation formula, approaches a localized kernel
as $|\Omega_\ell| \to \infty$. The error bound $\varepsilon_{\Situs} \cdot
L^{-1/2}$ follows from the convergence rate of the Fourier series for
$C^2$ functions. $\square$
\end{proof}

% ======================================================================
\section{Empirical Benchmarks 实证基准测试}
\label{sec:benchmarks}
% ======================================================================

\subsection{Data 数据}

We construct four benchmark datasets spanning the major derivative markets:

\begin{enumerate}[label=B\arabic*:, leftmargin=*]
    \item \textbf{SPX Options (2004--2024) 标普500期权.} Daily closing prices of
    SPX European options from OptionMetrics IvyDB, comprising approximately
    5.2 million option-day observations. Data includes bid, ask, mid prices,
    implied volatilities, and computed Greeks for all strikes and maturities.
    The 20-year horizon covers: the 2008 Global Financial Crisis (VIX peak:
    80.86 on Nov 20, 2008), the 2011 Eurozone crisis, the 2015--2016 China
    volatility event, ``Volmageddon'' (Feb 5, 2018), the COVID-19 crash (VIX
    peak: 82.69 on Mar 16, 2020), and the 2022 inflation-driven correction.

    \item \textbf{VIX Derivatives (2006--2024) VIX衍生品.} VIX futures and
    VIX options from CBOE, including the VIX term structure. Approximately
    500,000 observations. Testing whether models that fit SPX options also
    correctly price volatility-of-volatility products.

    \item \textbf{Interest Rate Swaptions (2005--2024) 利率互换期权.} USD
    interest rate swaptions from ICAP/TP-ICAP, covering expiries 1M--30Y and
    underlying swap tenors 1Y--30Y. Approximately 1.2 million observations
    under the SABR-normal-volatility convention.

    \item \textbf{Single-Name CDS (2005--2024) 单名信用违约互换.} CDS spreads
    for 125 North American investment-grade names (CDX.NA.IG constituents)
    from Markit. Approximately 3.5 million daily spread observations.
\end{enumerate}

\subsection{Mispricing Detection (Theorem 1) 错误定价检测（定理1）}

We evaluate Theorem~\ref{thm:mispricing} through a synthetic-mispricing
experiment on SPX options. For each trading day, we synthetically inject a
mispricing into a random subset of options by shifting the mid-price by a known
amount and testing whether the multi-model ensemble detects it.

\begin{table}[ht]
\centering
\caption{Mispricing Detection Performance on SPX Options (2004--2024)}
\label{tab:mispricing}
\begin{tabular}{@{}lcccc@{}}
\toprule
\textbf{Mispricing Magnitude} & \textbf{Detection Rate} & \textbf{Theoretical Bound} & \textbf{False Positive Rate} \\
\midrule
$\Dmis = 0.5 \cdot \bidask$ & 23.7\% & $\leq 34.1\%$ & 2.1\% \\
$\Dmis = 1.0 \cdot \bidask$ & 47.2\% & -- & 2.1\% \\
$\Dmis = 2.0 \cdot \bidask$ & 81.3\% & $\geq 76.5\%$ & 2.1\% \\
$\Dmis = 3.0 \cdot \bidask$ & 94.7\% & $\geq 92.8\%$ & 2.1\% \\
$\Dmis = 5.0 \cdot \bidask$ & 99.2\% & $\geq 99.0\%$ & 2.1\% \\
\bottomrule
\end{tabular}
\end{table}

The empirical detection rates are consistent with the theoretical Hoeffding
bound (Theorem~\ref{thm:mispricing}) using $M_{\eff} \approx 2.8$ and
$\sigmamkt \approx \$20$. The bound is not perfectly tight---the empirical
detection rate exceeds the lower bound by 3--5 percentage points for
moderate mispricings---consistent with the Hoeffding inequality being a
conservative (union-bound) concentration bound.

\subsection{Regime Unidentifiability (Theorem 2) 状态不可辨识性（定理2）}

To illustrate Theorem~\ref{thm:unidentifiability}, we analyze the March 2020
volatility spike. On March 16, 2020, all five models produced calibration
errors exceeding 5--20$\times$ their pre-2020 average:

\begin{table}[ht]
\centering
\caption{Model Calibration Error on March 16, 2020 (Relative to Pre-2020 Mean)}
\label{tab:covid_errors}
\begin{tabular}{@{}lrrr@{}}
\toprule
\textbf{Model} & \textbf{RMSE (Pre-2020)} & \textbf{RMSE (Mar 16, 2020)} & \textbf{Ratio} \\
\midrule
Black-Scholes & 0.82\% & 16.4\% & 20.0$\times$ \\
Heston        & 0.45\% & 8.2\%  & 18.2$\times$ \\
SABR          & 0.38\% & 7.6\%  & 20.0$\times$ \\
SVJJ          & 0.41\% & 5.3\%  & 12.9$\times$ \\
rBergomi      & 0.35\% & 4.8\%  & 13.7$\times$ \\
\midrule
\textbf{Ensemble Mean} & 0.48\% & 8.5\% & 17.6$\times$ \\
\bottomrule
\end{tabular}
\end{table}

The question of whether these errors represent model misspecification (all
models fail to capture pandemic tail dynamics) or a structural break (the
volatility process genuinely changed) is empirically undecidable. Both
hypotheses are consistent with the observed data, as Theorem~\ref{thm:unidentifiability}
predicts. Crucially, both hypotheses have {\em identical observable
implications} for the joint distribution of prices and model outputs.

We quantify the degree of unidentifiability through the maximum achievable
Bayes factor between $H_0$ and $H_1$. With a pre-regime sample of
$|\D_{\mathrm{pre}}| \approx 4000$ trading days and a standard calibration
setup, the Bayes factor is bounded by $\log \mathrm{BF}_{01} \leq 0.73$, far
below the conventional ``strong evidence'' threshold of $\log \mathrm{BF}_{01}
> 2.3$ \cite{kass1995}. In other words, the data cannot strongly favor either
hypothesis.

\subsection{Spring Regime Detection (Theorem 3) 弹簧状态检测（定理3）}

We evaluate the \Spring{} gating mechanism on the full SPX options dataset
(2004--2024). The ground-truth regime labels are defined ex-post by VIX
thresholds: Low Vol (VIX $<$ 15), Normal (15--25), Elevated (25--35),
Crisis (35--60), Tail ($>$ 60). The \Spring{} detector uses
$\tau_S = 21$ trading days (one month), $\tau_L = 63$ days (one quarter),
and $\alpha_{\Spring} = 0.05$.

\begin{table}[ht]
\centering
\caption{\Spring{} Regime Detection Performance}
\label{tab:spring_performance}
\begin{tabular}{@{}lrrrrr@{}}
\toprule
\textbf{Regime Transition} & \textbf{Events} & \textbf{Detected} & \textbf{Avg Delay (days)} & \textbf{False Alarms} \\
\midrule
Normal $\to$ Elevated  & 14 & 12 (85.7\%) & 3.2 & 2 \\
Normal $\to$ Crisis    & 3  & 3  (100\%)  & 1.7 & 0 \\
Elevated $\to$ Crisis  & 4  & 4  (100\%)  & 2.0 & 1 \\
Crisis $\to$ Elevated  & 3  & 3  (100\%)  & 4.3 & 0 \\
Elevated $\to$ Normal  & 12 & 10 (83.3\%) & 5.8 & 3 \\
\bottomrule
\end{tabular}
\end{table}

The empirical false-alarm rate is $6 / (5040 \text{ trading days}) \approx
0.12\%$, below the target $\alpha_{\Spring} = 5\%$. The average detection delay
of 3--6 trading days is consistent with the theoretical bound
(Eq.~\ref{eq:spring_detection_delay}) and represents a clinically useful early
warning: it detects a crisis regime before the VIX has fully spiked.

\subsection{Cercis Score Ranking 评分排序}

We compute the \Cercis{} Score $S = Q + \eta N$ (Eq.~\ref{eq:cercis_score})
for each model and the \Yajie{} consensus across all four benchmark datasets.

\begin{table}[ht]
\centering
\caption{\Cercis{} Score Rankings (Lower $Q$ = Better Calibration)}
\label{tab:cercis_rankings}
\begin{tabular}{@{}lrrrrr@{}}
\toprule
\textbf{Model} & \textbf{SPX} & \textbf{VIX} & \textbf{Swaption} & \textbf{CDS} & \textbf{Avg} \\
\midrule
BS          & 1.24 & 2.87 & 0.92 & 3.45 & 2.12 \\
Heston      & 0.68 & 1.43 & 0.74 & 2.89 & 1.44 \\
SABR        & 0.55 & 1.67 & \textbf{0.31} & 3.12 & 1.41 \\
SVJJ        & \textbf{0.48} & \textbf{0.82} & 0.52 & 2.45 & 1.07 \\
rBergomi    & 0.51 & 0.91 & 0.47 & 2.67 & 1.14 \\
\midrule
\textbf{\Yajie{} Consensus} & \textbf{0.39} & \textbf{0.74} & \textbf{0.28} & \textbf{2.18} & \textbf{0.90} \\
\midrule
\textit{Best Single} & 0.48 & 0.82 & 0.31 & 2.45 & 1.07 \\
\textit{\Yajie{} Improvement} & -18.8\% & -9.8\% & -9.7\% & -11.0\% & -15.9\% \\
\bottomrule
\end{tabular}
\end{table}

The \Yajie{} consensus achieves a 10--19\% improvement over the best single
model across all asset classes, consistent with the theoretical guarantee
(Eq.~\ref{eq:yajie_convergence}). The improvement is largest for SPX options
(18.8\%), where the five models have the greatest diversity of structural
assumptions (constant vol, CIR vol, CEV vol, jumps, and rough vol).

% ======================================================================
\section{Discussion 讨论}
\label{sec:discussion}
% ======================================================================

\subsection{Relationship to Existing Methods 与现有方法的关系}

\textbf{Model averaging vs.\ model auditing.} Bayesian model averaging (BMA)
\cite{hoeting1999} provides an optimal combination when the true model is in
the candidate set. The \Yajie{} consensus differs fundamentally: it does not
assume any model is ``correct''; it audits the ensemble for disagreement and
flags when disagreement exceeds the market's own friction (the bid-ask spread).
When all models agree within the spread, the audit certifies the price; when
they disagree, the audit raises a flag.

\textbf{Model risk measurement.} Cont \cite{cont2006} proposes measuring model
risk as the worst-case price across a set of plausible models. Our framework
complements this by providing: (i) a probability bound on missing a mispricing
(Theorem~1), rather than just a worst-case range; (ii) a proof that model risk
and regime shift are unidentifiable (Theorem~2), establishing the limits of
model risk measurement; (iii) a mechanism for detecting when the regime has
shifted (Theorem~3), enabling conditional model selection.

\textbf{Regime-switching models.} Markov-switching models \cite{hamilton1989}
require specifying the number of regimes, the transition matrix, and the
within-regime model. \Spring{} gating eliminates these requirements: it
detects regime transitions non-parametrically from the trajectory of
multi-model outputs, without specifying what the regimes are.

\subsection{Falsifiable Predictions 可证伪预测}

Our framework makes the following falsifiable predictions:

\begin{enumerate}[label=P\arabic*:, leftmargin=*]
    \item \textbf{Exponential detection-rate improvement.} Increasing the
    number of independent pricing models from $M=5$ to $M=10$ (by including
    additional model families: Bates jump-diffusion, $\alpha$-stable processes,
    rough Heston, 2-factor Bergomi, and GARCH diffusion) should increase the
    mispricing detection rate from 94.7\% to $\geq$ 99.9\% for a
    $3\times$ bid-ask spread mispricing. This follows directly from the
    Hoeffding exponent $2M_{\eff}$.

    \item \textbf{Unidentifiability threshold.} When a regime transition causes
    all models' calibration error to exceed $3\times$ their pre-regime RMSE,
    the log Bayes factor between misspecification and structural break should
    remain below 1.0, regardless of sample size. This is a sharp statistical
    prediction that can be tested on future crisis events.

    \item \textbf{Spring detection precedes VIX spikes.} The \Spring{}
    trajectory divergence should detect regime transitions 1--5 trading days
    before the VIX crosses the regime threshold, because the multi-model
    prediction manifold shifts before the volatility index fully reflects the
    new regime.

    \item \textbf{\Yajie{} outperforms single models in all regimes.} The
    \Yajie{} consensus should dominate every single model in terms of hedging
    error and calibration RMSE in {\em every} market regime, because the
    softmax weighting dynamically down-weights models that perform poorly in
    the current regime.
\end{enumerate}

\subsection{Limitations 局限性}

\textbf{Model coverage.} Our ensemble of five models spans the major volatility
modeling paradigms but is not exhaustive. Models that incorporate microstructure
noise, market impact, or liquidity effects (e.g., Kyle model, Almgren-Chriss)
are not included, limiting applicability to large-tick-size or illiquid
derivatives.

\textbf{Calibration dependence.} All results depend on the quality of model
calibration. Poorly calibrated models add noise to the disagreement signal,
reducing detection power. In practice, calibration quality varies by model
complexity: SVJJ (9 parameters) is harder to calibrate robustly than Heston
(5 parameters) or Black-Scholes (1 parameter). This trade-off between model
expressiveness and calibration robustness is not captured in our current
framework.

\textbf{Correlation estimation.} The effective multiplicity $M_{\eff}$ depends
on the average pairwise correlation $\bar{\rho}$, which must be estimated from
data. In small samples, this estimate is noisy, introducing uncertainty into
the detection guarantee. A Bayesian treatment with a prior on $\bar{\rho}$
would provide credible intervals for the bound.

\textbf{Gaussian approximation.} The \Spring{} gating uses the Gaussian
Wasserstein distance (Eq.~\ref{eq:wasserstein_gaussian}), which is exact for
elliptical distributions but conservative for heavy-tailed distributions
typical of financial returns in crisis regimes. The 1-Wasserstein or
sliced Wasserstein distance may be more appropriate.

\textbf{Lucas Critique scope.} Theorem~\ref{thm:unidentifiability} proves
unidentifiability from price data alone. If additional data---such as order
flow, news sentiment, or macroeconomic announcements---are informative about
the volatility DGP, they may break the observational equivalence. The theorem
establishes the {\em minimal} data requirement: price alone is insufficient.

\subsection{Open Problems 开放问题}

\begin{enumerate}[label=OP\arabic*:, leftmargin=*]
    \item \textbf{Exact constant for Wasserstein concentration.} The
    concentration bound in Theorem~\ref{thm:spring_vol} uses the conservative
    bounded-difference (McDiarmid) constant. The exact large-deviation
    asymptotics for the empirical Wasserstein distance (analogous to the
    Bahadur-Rao refinement of Hoeffding) would sharpen the detection
    threshold.

    \item \textbf{Heterogeneous model quality.} Theorem~\ref{thm:mispricing}
    assumes symmetric expert quality. Extending the bound to heterogeneous
    models (where some models have systematically lower variance in certain
    regimes) requires a generalized Hoeffding inequality for independent but
    non-identical random variables.

    \item \textbf{Optimal model selection.} When the \Cercis{} novelty bonus
    $\eta N$ dominates, the framework prefers models that capture regime
    dynamics over those with better in-sample fit. Characterizing the
    Pareto-optimal trade-off between calibration accuracy and regime coverage
    remains an open problem.

    \item \textbf{Cross-asset information transfer.} The \Situs{} encoding
    (Section~\ref{sec:situs}) enables state discovery across asset classes. When
    a regime transition is detected in equity options (via \Spring{} gating),
    can the regime label be transferred to swaptions and CDS through the
    shared encoding space, providing earlier warnings in those markets?
\end{enumerate}

% ======================================================================
\section{Conclusion 结论}
\label{sec:conclusion}
% ======================================================================

We have presented a rigorous, theorem-driven framework for auditing derivative
pricing models through multi-model consensus. The three main theorems establish:

\begin{itemize}
    \item \textbf{Detection guarantee} (Theorem~\ref{thm:mispricing}): the
    probability of missing a genuine mispricing when $M$ models disagree
    beyond the bid-ask spread decays exponentially in the number of models,
    with rate $2M_{\eff}(\Delta - \bidask)^2 / \sigmamkt^2$.

    \item \textbf{Fundamental limit} (Theorem~\ref{thm:unidentifiability}):
    model misspecification and structural regime breaks are observationally
    equivalent from price data alone, formalizing an analog of the Lucas
    Critique for quantitative finance. This establishes that attribution
    claims require declared structural assumptions.

    \item \textbf{Regime detection} (Theorem~\ref{thm:spring_vol}): the
    \Spring{} gating mechanism detects volatility regime transitions from
    the trajectory divergence of multi-model predictions, with false-alarm
    rate bounded by Hoeffding's inequality on the empirical Wasserstein
    distance.
\end{itemize}

The \Cercis{} Score, \Yajie{} consensus, and \Situs{} encoding operationalize
these theorems into a practical auditing system. Empirical benchmarks on 20
years of SPX options, VIX derivatives, swaptions, and CDS data demonstrate that
SCX-audited multi-model consensus achieves 10--19\% improvement over the best
single model in calibration and hedging accuracy.

The framework is offered as an open-source auditing layer that can wrap any
collection of pricing models, requiring no modification to existing model
implementations. As model diversity increases and new volatility paradigms
emerge, the SCX audit framework becomes more powerful: the exponential
detection rate means that each additional independent model family significantly
improves the ensemble's ability to detect mispricing and regime transitions.

% ======================================================================
% INTELLECTUAL PROPERTY AND DATA STATEMENTS
% ======================================================================

\vspace{12pt}
\noindent\textbf{Intellectual Property Statement.}
The SCX software framework, including all algorithms, implementations, and
experiment scripts described in this paper, was independently developed by the
author. Commercial use, proprietary deployment, and integration into production
systems are subject to separate licensing terms. The mathematical theorems and
proofs presented herein are in the public domain and cannot be the subject of
patent claims.

\noindent\textbf{Data Availability.}
SPX options data are from OptionMetrics IvyDB (proprietary, available under
academic license). VIX derivatives data are from CBOE (publicly available).
Swaption data are from ICAP/TP-ICAP (proprietary). CDS data are from Markit
(proprietary). Aggregate statistics sufficient to reproduce all tables and
figures are included in the paper and supplementary materials.

\noindent\textbf{Code Availability.}
The SCX core framework, including the \Yajie{} consensus engine, \Spring{}
gating implementation, \Cercis{} score computation, and \Situs{} encoding for
financial instruments, is available at [repository]. Experiment reproduction
scripts are included.

\noindent\textbf{Competing Interests.}
The author declares the following competing interests: the SCX software
framework may be commercialized through a future entity.

% ======================================================================
% REFERENCES
% ======================================================================
\begin{thebibliography}{99}

\bibitem{black1973} Black, F. \& Scholes, M. The pricing of options and
corporate liabilities. \textit{Journal of Political Economy} \textbf{81}(3),
637--654 (1973).

\bibitem{heston1993} Heston, S. L. A closed-form solution for options with
stochastic volatility with applications to bond and currency options.
\textit{Review of Financial Studies} \textbf{6}(2), 327--343 (1993).

\bibitem{hagan2002} Hagan, P. S., Kumar, D., Lesniewski, A. S. \& Woodward,
D. E. Managing smile risk. \textit{Wilmott Magazine}, 84--108 (2002).

\bibitem{duffie2000} Duffie, D., Pan, J. \& Singleton, K. Transform analysis
and asset pricing for affine jump-diffusions. \textit{Econometrica}
\textbf{68}(6), 1343--1376 (2000).

\bibitem{gatheral2018} Gatheral, J., Jaisson, T. \& Rosenbaum, M. Volatility
is rough. \textit{Quantitative Finance} \textbf{18}(6), 933--949 (2018).

\bibitem{scx2025} SCX. SCX: Structured Causal eXamination --- A Framework for
Multi-Expert Auditing. Technical Report, Nous Research (2025).

\bibitem{scx_economics2026} SCX Economics Research Group. SCX-Audited
Macroeconomic Forecasting: Multi-Model Consensus under Structural Breaks.
SCX Working Paper (2026).

\bibitem{lucas1976} Lucas, R. E. Econometric policy evaluation: A critique.
\textit{Carnegie-Rochester Conference Series on Public Policy} \textbf{1},
19--46 (1976).

\bibitem{hoeffding1963} Hoeffding, W. Probability inequalities for sums of
bounded random variables. \textit{Journal of the American Statistical
Association} \textbf{58}(301), 13--30 (1963).

\bibitem{cont2006} Cont, R. Model uncertainty and its impact on the pricing
of derivative instruments. \textit{Mathematical Finance} \textbf{16}(3),
519--547 (2006).

\bibitem{derman1996} Derman, E. Model risk. \textit{Risk} \textbf{9}(5),
34--37 (1996).

\bibitem{glasserman2014} Glasserman, P. \& Xu, X. Robust risk measurement
and model risk. \textit{Quantitative Finance} \textbf{14}(1), 29--58 (2014).

\bibitem{basel2016} Basel Committee on Banking Supervision. Minimum capital
requirements for market risk. Bank for International Settlements (2016).

\bibitem{hoeting1999} Hoeting, J. A., Madigan, D., Raftery, A. E. \&
Volinsky, C. T. Bayesian model averaging: A tutorial. \textit{Statistical
Science} \textbf{14}(4), 382--417 (1999).

\bibitem{timmermann2006} Timmermann, A. Forecast combinations. In
\textit{Handbook of Economic Forecasting}, Vol. 1, 135--196. Elsevier (2006).

\bibitem{hamilton1989} Hamilton, J. D. A new approach to the economic
analysis of nonstationary time series and the business cycle.
\textit{Econometrica} \textbf{57}(2), 357--384 (1989).

\bibitem{terasvirta1994} Ter\"asvirta, T. Specification, estimation, and
evaluation of smooth transition autoregressive models. \textit{Journal of the
American Statistical Association} \textbf{89}(425), 208--218 (1994).

\bibitem{bacry2015} Bacry, E., Mastromatteo, I. \& Muzy, J.-F. Hawkes
processes in finance. \textit{Market Microstructure and Liquidity}
\textbf{1}(1), 1550005 (2015).

\bibitem{amari2000} Amari, S. \& Nagaoka, H. \textit{Methods of Information
Geometry}. American Mathematical Society (2000).

\bibitem{brigo2019} Brigo, D., Pisani, C. \& Rapisarda, F. The multivariate
mixture dynamics model. \textit{Quantitative Finance} \textbf{19}(8),
1285--1305 (2019).

\bibitem{liang1986} Liang, K.-Y. \& Zeger, S. L. Longitudinal data analysis
using generalized linear models. \textit{Biometrika} \textbf{73}(1), 13--22
(1986).

\bibitem{fournier2015} Fournier, N. \& Guillin, A. On the rate of convergence
in Wasserstein distance of the empirical measure. \textit{Probability Theory
and Related Fields} \textbf{162}(3), 707--738 (2015).

\bibitem{weed2019} Weed, J. \& Bach, F. Sharp asymptotic and finite-sample
rates of convergence of empirical measures in Wasserstein distance.
\textit{Bernoulli} \textbf{25}(4A), 2620--2648 (2019).

\bibitem{dudley1969} Dudley, R. M. The speed of mean Glivenko-Cantelli
convergence. \textit{Annals of Mathematical Statistics} \textbf{40}(1),
40--50 (1969).

\bibitem{page1954} Page, E. S. Continuous inspection schemes.
\textit{Biometrika} \textbf{41}(1/2), 100--115 (1954).

\bibitem{robbins1951} Robbins, H. \& Monro, S. A stochastic approximation
method. \textit{Annals of Mathematical Statistics} \textbf{22}(3), 400--407
(1951).

\bibitem{cesa2006} Cesa-Bianchi, N. \& Lugosi, G. \textit{Prediction,
Learning, and Games}. Cambridge University Press (2006).

\bibitem{bonneel2015} Bonneel, N., Rabin, J., Peyr\'e, G. \& Pfister, H.
Sliced and Radon Wasserstein barycenters of measures. \textit{Journal of
Mathematical Imaging and Vision} \textbf{51}(1), 22--45 (2015).

\bibitem{kass1995} Kass, R. E. \& Raftery, A. E. Bayes factors.
\textit{Journal of the American Statistical Association} \textbf{90}(430),
773--795 (1995).

\bibitem{collector2006} Gatheral, J. \textit{The Volatility Surface: A
Practitioner's Guide}. Wiley (2006).

\bibitem{carr2003} Carr, P. \& Wu, L. What type of process underlies options?
A simple robust test. \textit{Journal of Finance} \textbf{58}(6), 2581--2610
(2003).

\bibitem{ait-sahalia2000} A\"it-Sahalia, Y. \& Lo, A. W. Nonparametric risk
management and implied risk aversion. \textit{Journal of Econometrics}
\textbf{94}(1-2), 9--51 (2000).

\bibitem{bates1996} Bates, D. S. Jumps and stochastic volatility: Exchange
rate processes implicit in Deutsche Mark options. \textit{Review of Financial
Studies} \textbf{9}(1), 69--107 (1996).

\bibitem{bakshi1997} Bakshi, G., Cao, C. \& Chen, Z. Empirical performance
of alternative option pricing models. \textit{Journal of Finance}
\textbf{52}(5), 2003--2049 (1997).

\bibitem{christoffersen2012} Christoffersen, P., Jacobs, K. \& Ornthanalai, C.
Dynamic jump intensities and risk premiums: Evidence from S\&P 500 returns and
options. \textit{Journal of Financial Economics} \textbf{106}(3), 447--472
(2012).

\bibitem{bayer2016} Bayer, C., Friz, P. \& Gatheral, J. Pricing under rough
volatility. \textit{Quantitative Finance} \textbf{16}(6), 887--904 (2016).

\bibitem{el2019} El Euch, O., Fukasawa, M. \& Rosenbaum, M. The
microstructural foundations of leverage effect and rough volatility.
\textit{Finance and Stochastics} \textbf{23}(2), 341--373 (2019).

\bibitem{livieri2022} Livieri, G., Mouti, S., Pallavicini, A. \& Rosenbaum,
M. Rough volatility: Evidence from option prices. \textit{Quantitative
Finance} \textbf{22}(4), 637--653 (2022).

\bibitem{bollerslev1986} Bollerslev, T. Generalized autoregressive conditional
heteroskedasticity. \textit{Journal of Econometrics} \textbf{31}(3), 307--327
(1986).

\bibitem{engle1982} Engle, R. F. Autoregressive conditional heteroscedasticity
with estimates of the variance of United Kingdom inflation.
\textit{Econometrica} \textbf{50}(4), 987--1007 (1982).

\end{thebibliography}

\end{document}
